\chapter{The \texorpdfstring{$\lambda$}{Lambda}-Calculus and Types}\label{ch:lambda}

\begin{center}\itshape
A function is a rule, a rule can be applied to anything, even to a rule;\\
and a proof of ``$A$ implies $B$'' is a program turning proofs of $A$ into proofs of $B$.\\
\upshape(the content of this chapter, in our words)
\end{center}

In 1928 Alonzo Church, a young assistant professor at Princeton, began to build a new foundation for logic in which the basic notion would not be the set but the function. To speak of functions without naming them, he wrote $\lambda x.\,x^2+1$ for ``the function which sends $x$ to $x^2+1$'', and he allowed any function to be applied to any argument, including itself. His papers of 1932 and 1933 set out the system. In 1935 two of his graduate students, Stephen Kleene and Barkley Rosser, proved that it was inconsistent: it could derive a contradiction.

What survived the disaster was its core, the rules for forming functions and applying them, without any logic: the \emph{$\lambda$-calculus}. It looked too poor to do anything, since even the predecessor function $n\mapsto n-1$ seemed out of reach, even to Church; Kleene found how to define it, and then so many other functions that Church came to believe, in 1934, that every effectively calculable function of the integers is definable in the $\lambda$-calculus. He proposed this to Kurt Gödel, who found the proposal ``thoroughly unsatisfactory''; Church published it anyway, in 1936, together with the first undecidability theorem. Turing's machines, a few months later, convinced Gödel, and Turing proved that they compute exactly Church's functions.

The $\lambda$-calculus then lived a second life. John McCarthy took its notation for the programming language LISP in 1958, and every functional language since descends from it. And the types, which Church added in 1940 to restore a logic, turned out to be logic itself: Haskell Curry had noticed in 1934, and William Howard explained in a manuscript of 1969, that the types of programs are the formulas of logic, and the programs are their proofs. Nicolaas de Bruijn built on this idea, from 1967, the first proof checker, Automath; its descendants have since checked the four-colour theorem and the proof that five-state busy beavers halt (Chapter~\ref{ch:turing}).

\begin{think}
The function $\lambda x.\,x\,x$ applies its argument to itself. What happens when you give it to itself? Can a type system, which forbids such things, still compute every function you care about? And what does it mean for a program to be a proof?
\end{think}

\section{Terms and reduction}\label{sec:terms}

\begin{definition}\label{def:lambdaterm}
Fix an infinite set of variables. The \emph{$\lambda$-terms} are generated by: every variable $x$ is a term; if $M,N$ are terms, so is the \emph{application} $MN$; if $M$ is a term and $x$ a variable, the \emph{abstraction} $\lambda x.M$ is a term. Application associates to the left ($MNP$ is $(MN)P$), the body of an abstraction extends as far right as possible, and $\lambda xy.M$ abbreviates $\lambda x.\lambda y.M$. An occurrence of $x$ in $M$ is \emph{bound} if it is inside some $\lambda x.\,\cdot$, and \emph{free} otherwise; $\mathrm{FV}(M)$ is the set of free variables, and $M$ is \emph{closed} if it is empty. Terms differing only by the names of their bound variables, such as $\lambda x.x$ and $\lambda y.y$, are identified (\emph{$\alpha$-conversion}). The \emph{substitution} $M[x:=N]$ replaces the free occurrences of $x$ in $M$ by $N$, after renaming the bound variables of $M$ so that no free variable of $N$ gets captured.
\end{definition}

We always assume, renaming if necessary, that the bound variables of the terms in a statement are distinct from each other and from the free ones. The one computation rule is to apply a function.

\begin{definition}\label{def:beta}
A \emph{redex} is a term $(\lambda x.M)N$, and its \emph{contractum} is $M[x:=N]$. We write $M\to N$ (one step of \emph{$\beta$-reduction}) if $N$ is obtained from $M$ by replacing one redex occurring in $M$ by its contractum; $\twoheadrightarrow$ is the reflexive and transitive closure of $\to$, and $=_\beta$ the equivalence relation it generates. A term with no redex is in \emph{normal form}; $M$ \emph{has} a normal form $N$ if $M\twoheadrightarrow N$ with $N$ normal. A term is \emph{strongly normalising} if it has no infinite sequence of reductions $M\to M_1\to M_2\to\cdots$.
\end{definition}

A few closed terms deserve names: $\mathsf I=\lambda x.x$, $\mathsf K=\lambda xy.x$, $\mathsf S=\lambda xyz.xz(yz)$, and
\[
\Omega=(\lambda x.xx)(\lambda x.xx)\to(\lambda x.xx)(\lambda x.xx)=\Omega .
\]
So $\Omega$ has no normal form: its only reduct is itself. The term $\mathsf K\mathsf I\Omega$ shows that the order of reductions matters: $\mathsf K\mathsf I\Omega\to(\lambda y.\mathsf I)\Omega\to\mathsf I$, but reducing $\Omega$ forever is also possible. A term may reach normal form along one path and not along another; Figure~\ref{fig:lambdagraph} shows the more pleasant situation where all paths meet.

\begin{figure}[ht]
\centering
\fig{ch05-graph}
\caption{All the reductions of $(\lambda x.xx)(\mathsf I\mathsf I)$, one arrow per redex (the term $\mathsf I(\mathsf I\mathsf I)$ has two redexes with the same contractum). The paths diverge at once, one duplicating the argument and one simplifying it first, and all meet again at the normal form $\mathsf I$. Theorem~\ref{thm:churchrosser} says that paths can always be made to meet.}
\label{fig:lambdagraph}
\end{figure}

\section{Confluence}\label{sec:confluence}

\begin{theorem}[Church, Rosser, 1936]\label{thm:churchrosser}
If $M\twoheadrightarrow N_1$ and $M\twoheadrightarrow N_2$, there is $P$ with $N_1\twoheadrightarrow P$ and $N_2\twoheadrightarrow P$.
\end{theorem}

One step does not have this property (Exercise~\ref{ex:diamond}): reducing a redex may duplicate another one, which then needs several steps. The proof, due to William Tait and Per Martin-Löf, with a simplification by Masako Takahashi, uses steps that contract many redexes at once.

\begin{definition}\label{def:parallel}
\emph{Parallel reduction} $\Rightarrow$ is the smallest relation such that: $x\Rightarrow x$; if $M\Rightarrow M'$ then $\lambda x.M\Rightarrow\lambda x.M'$; if $M\Rightarrow M'$ and $N\Rightarrow N'$ then $MN\Rightarrow M'N'$ and $(\lambda x.M)N\Rightarrow M'[x:=N']$.
\end{definition}

By induction on terms, $M\Rightarrow M$; hence a single step $M\to N$ is a parallel step, and a parallel step contracts redexes one at a time, so $M\Rightarrow N$ implies $M\twoheadrightarrow N$. Consequently $\twoheadrightarrow$ is also the reflexive transitive closure of $\Rightarrow$. We need the following identity, proved by induction on $M$: if $x\ne y$ and $y\notin\mathrm{FV}(L)$, then
\begin{equation}\label{eq:substlemma}
M[y:=N][x:=L]=M[x:=L][y:=N[x:=L]].
\end{equation}
(Both sides replace $x$ by $L$ and $y$ by $N$ with its own $x$ replaced by $L$; the only nontrivial case is $M=y$, where both sides are $N[x:=L]$.)

\begin{lemma}\label{lem:parsubst}
If $M\Rightarrow M'$ and $N\Rightarrow N'$, then $M[x:=N]\Rightarrow M'[x:=N']$.
\end{lemma}

\begin{proof}
Induction on the derivation of $M\Rightarrow M'$. If $M=x$, both sides are $N\Rightarrow N'$; if $M=y\ne x$, both are $y$. For $\lambda y.P\Rightarrow\lambda y.P'$ and $PQ\Rightarrow P'Q'$, apply the induction hypothesis to the parts. For $(\lambda y.P)Q\Rightarrow P'[y:=Q']$, with $y\notin\mathrm{FV}(N)$ by our convention,
\[
\bigl((\lambda y.P)Q\bigr)[x:=N]=(\lambda y.P[x:=N])\,Q[x:=N]
\]
\[
\Rightarrow P'[x:=N'][y:=Q'[x:=N']]=P'[y:=Q'][x:=N'],
\]
by the last rule with the induction hypotheses, and by \eqref{eq:substlemma}.
\end{proof}

\begin{definition}\label{def:development}
The \emph{complete development} $M^*$ of $M$ contracts all its redexes at once: $x^*=x$, $(\lambda x.M)^*=\lambda x.M^*$, $\bigl((\lambda x.M)N\bigr)^*=M^*[x:=N^*]$, and $(MN)^*=M^*N^*$ if $M$ is not an abstraction.
\end{definition}

\begin{lemma}[Takahashi]\label{lem:takahashi}
If $M\Rightarrow N$, then $N\Rightarrow M^*$.
\end{lemma}

\begin{proof}
Induction on the derivation of $M\Rightarrow N$. The cases $x\Rightarrow x$ and $\lambda x.P\Rightarrow\lambda x.P'$ are immediate. If $M=(\lambda x.P)Q\Rightarrow P'[x:=Q']=N$, the induction hypothesis gives $P'\Rightarrow P^*$ and $Q'\Rightarrow Q^*$, and Lemma~\ref{lem:parsubst} gives $N\Rightarrow P^*[x:=Q^*]=M^*$. If $M=PQ\Rightarrow P'Q'=N$ with $P$ not an abstraction, then $N\Rightarrow P^*Q^*=M^*$ by induction. Finally, if $M=PQ\Rightarrow P'Q'$ with $P=\lambda x.R$, then $P'=\lambda x.R'$ with $R\Rightarrow R'$, and $N=(\lambda x.R')Q'\Rightarrow R^*[x:=Q^*]=M^*$ by the last rule of Definition~\ref{def:parallel} and the induction hypotheses $R'\Rightarrow R^*$, $Q'\Rightarrow Q^*$.
\end{proof}

\begin{proof}[Proof of Theorem~\ref{thm:churchrosser}]
By Lemma~\ref{lem:takahashi}, if $M\Rightarrow N_1$ and $M\Rightarrow N_2$ then $N_1\Rightarrow M^*$ and $N_2\Rightarrow M^*$: parallel reduction has the \emph{diamond property}. Now $M\twoheadrightarrow N_1$ and $M\twoheadrightarrow N_2$ are sequences of parallel steps, of lengths $m$ and $n$. Fill the $m\times n$ grid of Figure~\ref{fig:diamonds} with diamonds, one at a time, from the corner $M$: each square, given its two upper sides, is completed by the diamond property. The last corner $P$ is reached from $N_1$ and from $N_2$ by parallel steps, hence by $\twoheadrightarrow$.
\end{proof}

\begin{figure}[ht]
\centering
\fig{ch05-diamonds}
\caption{The proof of the Church--Rosser theorem: each arrow is a parallel step. The given reductions from $M$ (blue) are completed by diamonds of parallel reduction (orange), square by square, until both reach $P$.}
\label{fig:diamonds}
\end{figure}

\begin{corollary}\label{cor:nf}
\textup{(a)} $M=_\beta N$ if and only if $M$ and $N$ have a common reduct. \textup{(b)} A term has at most one normal form, and if $M=_\beta N$ with $N$ normal then $M\twoheadrightarrow N$. \textup{(c)} The $\lambda$-calculus is consistent: distinct normal forms, such as $\mathsf K$ and $\mathsf I$, are not $\beta$-equal.
\end{corollary}

\begin{proof}
(a) Having a common reduct is reflexive, symmetric, contains $\to$, and is transitive: if $M,N$ reduce to $P$ and $N,L$ to $Q$, then $P,Q$ have a common reduct $R$ by the theorem applied to $N$, and $M,L\twoheadrightarrow R$. So it is $=_\beta$. (b) and (c): a common reduct of two normal forms is equal to both.
\end{proof}

Which path reaches the normal form? Always contracting the leftmost redex (the one whose $\lambda$ is leftmost) does, if any path does: this is the \emph{normalisation theorem} of Curry and Feys (1958), which we quote. In $\mathsf K\mathsf I\Omega$ the leftmost strategy throws $\Omega$ away before touching it. Programming languages that evaluate arguments before calling functions do not have this guarantee, and lazy languages such as Haskell are built on it.

\section{Programming with \texorpdfstring{$\lambda$}{lambda}}\label{sec:church}

The $\lambda$-calculus has no numbers, so we make them, by what they do: the number $n$ is the act of iterating a function $n$ times.

\begin{definition}\label{def:numerals}
The \emph{Church numerals} are $\overline n=\lambda fx.f^n(x)$, where $f^0(x)=x$ and $f^{n+1}(x)=f(f^n(x))$. A function $\varphi:\mathbb N^k\to\mathbb N$ is \emph{$\lambda$-definable} if there is a closed term $F$ with $F\,\overline{n_1}\cdots\overline{n_k}=_\beta\overline{\varphi(n_1,\dots,n_k)}$ for all $n_1,\dots,n_k$.
\end{definition}

\begin{proposition}\label{prop:arith}
Let $\mathsf{succ}=\lambda nfx.f(nfx)$, $\mathsf{plus}=\lambda mnfx.mf(nfx)$ and $\mathsf{mult}=\lambda mnf.m(nf)$. Then $\mathsf{succ}\,\overline n\twoheadrightarrow\overline{n+1}$, $\mathsf{plus}\,\overline m\,\overline n\twoheadrightarrow\overline{m+n}$ and $\mathsf{mult}\,\overline m\,\overline n\twoheadrightarrow\overline{mn}$.
\end{proposition}

\begin{proof}
First, $\overline n\,F\,X\twoheadrightarrow F^n(X)$ for all terms $F,X$, by two steps. Then $\mathsf{succ}\,\overline n\twoheadrightarrow\lambda fx.f(\overline nfx)\twoheadrightarrow\lambda fx.f(f^n(x))=\overline{n+1}$, and $\mathsf{plus}\,\overline m\,\overline n\twoheadrightarrow\lambda fx.\overline mf(\overline nfx)\twoheadrightarrow\lambda fx.f^m(f^n(x))=\overline{m+n}$. For the product, $\overline n f\twoheadrightarrow\lambda x.f^n(x)$, and iterating this function $m$ times gives $\lambda x.f^{mn}(x)$: by induction on $m$, $(\lambda x.f^n(x))^m(y)\twoheadrightarrow f^{mn}(y)$, so $\mathsf{mult}\,\overline m\,\overline n\twoheadrightarrow\lambda f.\overline m(\overline nf)\twoheadrightarrow\lambda fy.f^{mn}(y)=\overline{mn}$.
\end{proof}

Exponentiation is even simpler: $\overline n\,\overline m\twoheadrightarrow\overline{m^n}$ (Exercise~\ref{ex:exp}). For choices we encode truth values by what they do, choose: $\mathsf{true}=\lambda xy.x$, $\mathsf{false}=\lambda xy.y$, so that $\mathsf{true}\,M\,N\twoheadrightarrow M$ and $\mathsf{false}\,M\,N\twoheadrightarrow N$, and ``if $B$ then $M$ else $N$'' is just $BMN$. Then $\mathsf{iszero}=\lambda n.n(\lambda z.\mathsf{false})\mathsf{true}$ satisfies $\mathsf{iszero}\,\overline0\twoheadrightarrow\mathsf{true}$ and $\mathsf{iszero}\,\overline{n+1}\twoheadrightarrow(\lambda z.\mathsf{false})(\cdots)\to\mathsf{false}$. Pairs are $\langle M,N\rangle=\lambda z.zMN$, with $\mathsf{fst}=\lambda p.p\,\mathsf{true}$ and $\mathsf{snd}=\lambda p.p\,\mathsf{false}$, so that $\mathsf{fst}\langle M,N\rangle\twoheadrightarrow M$ and $\mathsf{snd}\langle M,N\rangle\twoheadrightarrow N$.

The predecessor is the function that stumped Church: a numeral applies $f$ some number of times, and there seems to be no way to apply it one time fewer. Kleene's trick is to iterate on pairs, keeping one step of history.

\begin{proposition}[Kleene]\label{prop:pred}
Let
\[
\mathsf{step}=\lambda p.\langle\mathsf{snd}\,p,\ \mathsf{succ}(\mathsf{snd}\,p)\rangle,\qquad\mathsf{pred}=\lambda n.\mathsf{fst}\bigl(n\,\mathsf{step}\,\langle\overline0,\overline0\rangle\bigr).
\]
Then $\mathsf{pred}\,\overline0\twoheadrightarrow\overline0$ and $\mathsf{pred}\,\overline{n+1}\twoheadrightarrow\overline n$.
\end{proposition}

\begin{proof}
Let $p_0=\langle\overline0,\overline0\rangle$ and $p_{k+1}=\mathsf{step}\,p_k$. By induction, $\mathsf{fst}\,p_{k+1}=_\beta\overline k$ and $\mathsf{snd}\,p_k=_\beta\overline k$: indeed $p_{k+1}\twoheadrightarrow\langle\mathsf{snd}\,p_k,\mathsf{succ}(\mathsf{snd}\,p_k)\rangle$. Since $\overline n\,\mathsf{step}\,p_0\twoheadrightarrow\mathsf{step}^n(p_0)=p_n$, we get $\mathsf{pred}\,\overline{n+1}=_\beta\mathsf{fst}\,p_{n+1}=_\beta\overline n$ and $\mathsf{pred}\,\overline0=_\beta\mathsf{fst}\,p_0=_\beta\overline0$, and by Corollary~\ref{cor:nf}(b) the reductions $\twoheadrightarrow$ hold, as numerals are normal.
\end{proof}

Recursion needs a function to call itself, and the $\lambda$-calculus has no names. It does not need them.

\begin{theorem}[fixed points]\label{thm:fixedpoint}
Let
\[
\mathsf Y=\lambda f.(\lambda x.f(xx))(\lambda x.f(xx)).
\]
For every term $F$, $\mathsf YF=_\beta F(\mathsf YF)$.
\end{theorem}

\begin{proof}
Let $W=(\lambda x.F(xx))(\lambda x.F(xx))$. Then $\mathsf YF\to W\to F(W)$, and $F(\mathsf YF)\to F(W)$.
\end{proof}

This is the diagonal argument once more (Chapter~\ref{ch:turing}), now as a construction instead of a contradiction: $\lambda x.F(xx)$ applied to itself produces $F$ of itself. To define a function by a recursive equation $h=G\,h$, take $h=\mathsf YG$. Turing found a variant $\Theta$ with $\Theta F\twoheadrightarrow F(\Theta F)$ (Exercise~\ref{ex:turingfix}).

\begin{theorem}[Kleene, 1936; Church, 1936; Turing, 1937]\label{thm:lambdadefinable}
A total function $\mathbb N^k\to\mathbb N$ is $\lambda$-definable if and only if it is computable.
\end{theorem}

\begin{proof}
If $F$ $\lambda$-defines $\varphi$, then $F\,\overline{n_1}\cdots\overline{n_k}\twoheadrightarrow\overline{\varphi(\bar n)}$ by Corollary~\ref{cor:nf}(b). An algorithm computes $\varphi(\bar n)$ by exploring all reduction paths of $F\,\overline{n_1}\cdots\overline{n_k}$ breadth first (each term has finitely many redexes) until it meets a numeral; this numeral is $\overline{\varphi(\bar n)}$, the unique normal form.

Conversely, by Kleene's normal form theorem (Theorem~\ref{thm:kleenenf}), a total computable $\varphi$ is $\varphi(\bar n)=V(\mu t\,[T(\bar n,t)=0])$ with $T,V$ primitive recursive, and for every $\bar n$ there is such a $t$. So it suffices to $\lambda$-define the primitive recursive functions, closed under composition and recursion, and this one minimisation.

\emph{Basic functions and composition.} Zero is $\lambda\bar x.\overline0$, successor $\mathsf{succ}$, projections $\lambda x_1\cdots x_k.x_i$; if $H,G_1,\dots,G_j$ define $h,g_1,\dots,g_j$, then $\lambda\bar x.H(G_1\bar x)\cdots(G_j\bar x)$ defines their composition.

\emph{Primitive recursion.} Let $\varphi(0,\bar x)=g(\bar x)$ and $\varphi(n+1,\bar x)=h(n,\varphi(n,\bar x),\bar x)$, with $g,h$ defined by $G,H$. Iterate on pairs as for the predecessor: let
\[
\mathsf{st}=\lambda p.\langle\mathsf{succ}(\mathsf{fst}\,p),\ H(\mathsf{fst}\,p)(\mathsf{snd}\,p)\bar x\rangle,\qquad\Phi=\lambda n\bar x.\mathsf{snd}\bigl(n\,\mathsf{st}\,\langle\overline0,G\bar x\rangle\bigr).
\] By induction on $n$, $\mathsf{st}^n\langle\overline0,G\bar x\rangle=_\beta\langle\overline n,\overline{\varphi(n,\bar x)}\rangle$ for numerals $\bar x$, so $\Phi$ defines $\varphi$.

\emph{Minimisation.} Let $T$ be defined by $T'$, and let
\[
\mathsf{Z}=\lambda h t.\,\mathsf{iszero}(T'\bar n\,t)\;t\;\bigl(h(\mathsf{succ}\,t)\bigr),
\]
a term with the numerals $\bar n$ as parameters. By Theorem~\ref{thm:fixedpoint}, $W=\mathsf Y\mathsf Z$ satisfies
\[
W\,\overline t=_\beta\mathsf Z\,W\,\overline t=_\beta\mathsf{iszero}(T'\bar n\,\overline t)\;\overline t\;(W\,\overline{t+1}),
\]
which is $=_\beta\overline t$ if $T(\bar n,t)=0$ and $=_\beta W\,\overline{t+1}$ otherwise. If $t_0$ is the least zero, then $W\,\overline0=_\beta W\,\overline1=_\beta\cdots=_\beta W\,\overline{t_0}=_\beta\overline{t_0}$. So $\lambda\bar n.\,V'\bigl(\mathsf Y(\mathsf Z)\,\overline0\bigr)$, with $\bar n$ abstracted, defines $\varphi$.
\end{proof}

The theorem extends to partial functions, with ``undefined'' represented by terms without normal form, a refinement we do not need. Church's 1936 paper also proved, by a diagonal argument like Turing's, that no algorithm decides whether a $\lambda$-term has a normal form, or whether two terms are $\beta$-equal; we quote it. Theorem~\ref{thm:lambdadefinable}, with Theorem~\ref{thm:kleenenf}, gives one more proof that the computable functions do not depend on the model.

\section{Simple types}\label{sec:types}

Self-application, which makes $\Omega$ loop and $\mathsf Y$ work, is exactly what types forbid.

\begin{definition}\label{def:types}
The \emph{simple types} are built from \emph{atomic types} $o,p,q,\dots$ by $A\to B$, with $A\to B\to C$ meaning $A\to(B\to C)$. A \emph{context} $\Gamma$ is a finite list $x_1:A_1,\dots,x_n:A_n$ of distinct variables with types. The \emph{typing judgement} $\Gamma\vdash M:A$ is derived by the rules
\[
\frac{x:A\in\Gamma}{\Gamma\vdash x:A}\qquad
\frac{\Gamma,x:A\vdash M:B}{\Gamma\vdash\lambda x.M:A\to B}\qquad
\frac{\Gamma\vdash M:A\to B\quad\Gamma\vdash N:A}{\Gamma\vdash MN:B}
\]
A term is \emph{typable} if $\Gamma\vdash M:A$ for some $\Gamma$ and $A$.
\end{definition}

For example $\vdash\mathsf I:A\to A$ and $\vdash\mathsf K:A\to B\to A$ for all types $A,B$, and every numeral has type $N_A=(A\to A)\to A\to A$. The term $\lambda x.xx$ is not typable: $x$ would need a type $A$ equal to $A\to B$, which is longer (Exercise~\ref{ex:untypable}). So neither $\Omega$ nor $\mathsf Y$ is typable, and the next theorem explains why they could not be.

The rules are syntax-directed: a derivation of $\Gamma\vdash\lambda x.M:C$ ends with the abstraction rule, so $C=A\to B$ with $\Gamma,x:A\vdash M:B$; a derivation of $\Gamma\vdash MN:B$ ends with the application rule. By induction on derivations one shows that adding variables to the context preserves judgements, and the \emph{substitution property}: if $\Gamma,x:A\vdash M:B$ and $\Gamma\vdash N:A$, then $\Gamma\vdash M[x:=N]:B$ (at each occurrence of $x$, use the derivation of $N$).

\begin{proposition}[subject reduction]\label{prop:subjred}
If $\Gamma\vdash M:A$ and $M\to M'$, then $\Gamma\vdash M':A$.
\end{proposition}

\begin{proof}
Induction on the position of the contracted redex. If $M$ itself is the redex $(\lambda x.P)Q$, then $\Gamma\vdash\lambda x.P:B\to A$ and $\Gamma\vdash Q:B$ for some $B$, so $\Gamma,x:B\vdash P:A$, and $\Gamma\vdash P[x:=Q]:A$ by the substitution property. If the redex is inside $P$ or $Q$ in $M=PQ$, or inside the body of an abstraction, the induction hypothesis applies to that part, with the same types.
\end{proof}

\begin{theorem}[Tait, 1967]\label{thm:sn}
Every typable term is strongly normalising.
\end{theorem}

The direct approach, induction on the term, fails: $MN$ may be strongly normalising in its parts, but contracting $(\lambda x.P)N$ creates $P[x:=N]$, which can be larger and contain new redexes. Tait's idea is to prove a stronger statement, defined by induction on types, which is preserved by application. Let $\mathrm{SN}$ be the set of strongly normalising terms; for $M\in\mathrm{SN}$, let $\nu(M)$ be the length of its longest reduction sequence, which is finite by König's lemma, since each term has finitely many one-step reducts. A term is \emph{neutral} if it is not an abstraction.

\begin{definition}\label{def:reducible}
The sets $\mathrm{Red}_A$ of \emph{reducible} terms of type $A$ are: $\mathrm{Red}_p=\mathrm{SN}$ for an atomic type $p$, and $\mathrm{Red}_{A\to B}=\{M:MN\in\mathrm{Red}_B\text{ for every }N\in\mathrm{Red}_A\}$.
\end{definition}

\begin{lemma}\label{lem:cr123}
For every type $A$:
\textup{(CR1)} $\mathrm{Red}_A\subseteq\mathrm{SN}$;
\textup{(CR2)} if $M\in\mathrm{Red}_A$ and $M\to M'$, then $M'\in\mathrm{Red}_A$;
\textup{(CR3)} if $M$ is neutral and every $M'$ with $M\to M'$ is in $\mathrm{Red}_A$, then $M\in\mathrm{Red}_A$.
In particular every variable is in $\mathrm{Red}_A$, by \textup{(CR3)}, since it has no reducts.
\end{lemma}

\begin{proof}
Induction on $A$. For an atomic type, (CR1) is the definition, (CR2) holds because the reducts of a strongly normalising term are strongly normalising, and (CR3) because an infinite reduction from $M$ would begin with a step to some $M'\in\mathrm{SN}$.

Let $A=B\to C$. (CR1) Let $M\in\mathrm{Red}_A$ and $x$ a variable; $x\in\mathrm{Red}_B$ by (CR3) for $B$, so $Mx\in\mathrm{Red}_C\subseteq\mathrm{SN}$ by (CR1) for $C$. An infinite reduction from $M$ would give one from $Mx$, so $M\in\mathrm{SN}$. (CR2) If $M\to M'$ and $N\in\mathrm{Red}_B$, then $MN\in\mathrm{Red}_C$ and $MN\to M'N$, so $M'N\in\mathrm{Red}_C$ by (CR2) for $C$. (CR3) Let $M$ be neutral with all its reducts in $\mathrm{Red}_A$, and $N\in\mathrm{Red}_B$; we show $MN\in\mathrm{Red}_C$ by induction on $\nu(N)$, finite by (CR1) for $B$. The term $MN$ is neutral, and since $M$ is not an abstraction, $MN$ is not a redex: its reducts are $M'N$ with $M\to M'$, in $\mathrm{Red}_C$ because $M'\in\mathrm{Red}_A$, and $MN'$ with $N\to N'$, in $\mathrm{Red}_C$ by the induction hypothesis since $N'\in\mathrm{Red}_B$ by (CR2) and $\nu(N')<\nu(N)$. By (CR3) for $C$, $MN\in\mathrm{Red}_C$.
\end{proof}

\begin{lemma}\label{lem:abstraction}
If $M[x:=N]\in\mathrm{Red}_B$ for every $N\in\mathrm{Red}_A$, then $\lambda x.M\in\mathrm{Red}_{A\to B}$.
\end{lemma}

\begin{proof}
Let $N\in\mathrm{Red}_A$. Then $M=M[x:=x]\in\mathrm{Red}_B\subseteq\mathrm{SN}$ and $N\in\mathrm{SN}$; we show $(\lambda x.M)N\in\mathrm{Red}_B$ by induction on $\nu(M)+\nu(N)$, using (CR3) for this neutral term. Its reducts are: $M[x:=N]$, in $\mathrm{Red}_B$ by hypothesis; $(\lambda x.M')N$ with $M\to M'$, where $M'$ satisfies the hypothesis of the lemma (each $M'[x:=N'']$ is a reduct of $M[x:=N'']$, use (CR2)) and $\nu(M')<\nu(M)$; and $(\lambda x.M)N'$ with $N\to N'$, where $N'\in\mathrm{Red}_A$ and $\nu(N')<\nu(N)$. The last two are in $\mathrm{Red}_B$ by the induction hypothesis.
\end{proof}

\begin{proof}[Proof of Theorem~\ref{thm:sn}]
We show by induction on derivations: if $x_1:A_1,\dots,x_n:A_n\vdash M:A$ and $N_i\in\mathrm{Red}_{A_i}$ for all $i$, then $M[\bar x:=\bar N]\in\mathrm{Red}_A$ (simultaneous substitution). For a variable this is the hypothesis. For $M=PQ$, $P[\bar x:=\bar N]\in\mathrm{Red}_{B\to A}$ and $Q[\bar x:=\bar N]\in\mathrm{Red}_B$ by induction, so their application is in $\mathrm{Red}_A$ by definition. For $M=\lambda y.P$ with $A=B\to C$, the induction hypothesis for $P$, with the context extended by $y:B$, gives $P[\bar x:=\bar N][y:=N]\in\mathrm{Red}_C$ for every $N\in\mathrm{Red}_B$, and Lemma~\ref{lem:abstraction} gives $\lambda y.P[\bar x:=\bar N]\in\mathrm{Red}_{B\to C}$. Finally, take $N_i=x_i$, which are reducible by Lemma~\ref{lem:cr123}: then $M\in\mathrm{Red}_A\subseteq\mathrm{SN}$.
\end{proof}

So typed programs always terminate, whatever the order of evaluation: a guarantee, and a limitation. By the diagonal argument, no programming language in which all programs terminate can compute all total computable functions (otherwise one would evaluate all its programs on their own codes and add one). The simply typed $\lambda$-calculus is very weak: Helmut Schwichtenberg showed in 1976 that on numerals of type $N_o$ it defines only the \emph{extended polynomials}, built from constants, $+$, $\times$ and a test for zero, a result we quote. Stronger type systems, such as Gödel's system T (1958) or Girard's system F (1972), keep termination and compute vastly more; the dividing line, again, is drawn by Gödel's incompleteness theorem.

\section{Proofs as programs}\label{sec:curryhoward}

Look at the typing rules with the terms erased:
\[
\frac{A\in\Gamma}{\Gamma\vdash A}\qquad
\frac{\Gamma,A\vdash B}{\Gamma\vdash A\to B}\qquad
\frac{\Gamma\vdash A\to B\quad\Gamma\vdash A}{\Gamma\vdash B}
\]
These are the rules of natural deduction for implication, as Gerhard Gentzen formulated them in 1935: an assumption may be used; to prove $A\to B$, assume $A$ and prove $B$, discharging the assumption; from $A\to B$ and $A$, conclude $B$ (\emph{modus ponens}). Formulas built from propositional variables by $\to$ and provable by these rules form \emph{minimal implicational logic}.

\begin{proposition}[Curry, 1934; Howard, 1969]\label{prop:curryhoward}
A formula $A$ is provable from assumptions $A_1,\dots,A_n$ if and only if there is a term $M$ with $x_1:A_1,\dots,x_n:A_n\vdash M:A$. More precisely, typed terms are notations for natural deduction proofs, and $\beta$-reduction is the simplification of proofs.
\end{proposition}

\begin{proof}
Erasing the terms of a typing derivation gives a proof, rule by rule. Conversely, a proof is turned into a derivation by induction: the use of the assumption $A_i$ becomes the variable $x_i$; the discharge of an assumption $A$ in a proof of $B$, given by a term $M$ with a new variable $x:A$, becomes $\lambda x.M:A\to B$; modus ponens becomes application. A redex $(\lambda x.M)N$ is a proof that proves $A\to B$ and immediately applies it to a proof $N$ of $A$; its contractum $M[x:=N]$ is the direct proof, where each use of the assumption $A$ is replaced by its proof $N$.
\end{proof}

Under this correspondence, Theorem~\ref{thm:sn} says that every proof can be put in normal form, a proof with no detour, and normal proofs have a rigid shape.

\begin{lemma}\label{lem:normalshape}
A term in normal form is either an abstraction $\lambda x.M$ or of the form $yM_1\cdots M_k$ ($k\ge0$) with $y$ a variable and $M_1,\dots,M_k$ normal. In particular a closed normal typed term has an arrow type, and no atomic type $p$ is provable without assumptions.
\end{lemma}

\begin{proof}
A term not an abstraction is $PM_1\cdots M_k$ with $P$ a variable or an abstraction; if $P$ were an abstraction and $k\ge1$, $PM_1$ would be a redex. If the term is closed, $y$ would be free, so it is an abstraction. If $\vdash M:p$, the normal form of $M$, which exists by Theorem~\ref{thm:sn} and has type $p$ by Proposition~\ref{prop:subjred}, would be a closed normal term of atomic type.
\end{proof}

Normal forms make questions of provability into finite searches. The formula $((p\to q)\to p)\to p$, \emph{Peirce's law}, is a tautology of classical logic, but it has no proof in minimal logic, because no closed term has this type (Exercise~\ref{ex:peirce}). What can be proved with implication alone is \emph{intuitionistic}: a proof must construct. Extended to the other connectives (pairs for $\wedge$, tagged unions for $\vee$, functions into the empty type for $\neg$) and to quantifiers (dependent types), the correspondence identifies the whole of constructive mathematics with a typed programming language. Classical logic corresponds to programs that capture their own continuation, as Timothy Griffin discovered in 1990: Peirce's law is the type of the operator \texttt{call/cc} of the language Scheme.

Two practical consequences close the chapter. First, types can be computed: J.~Roger Hindley showed in 1969 that every typable term has a \emph{principal type}, from which all its types are obtained by substitution, and that it is computed by unifying type equations; Robin Milner rediscovered the algorithm in 1978 for the language ML, where programmers rarely write types because the compiler infers them. We quote the theorem; the notebook implements the algorithm. Second, proofs can be checked by machines. A proof assistant is a type checker: a theorem is a type, a proof is a term, and checking the proof means checking that the term has that type. Automath (de Bruijn, 1967) was used by L.~S.~van Benthem Jutting to verify, in 1977, all of Landau's textbook \emph{Grundlagen der Analysis}; its descendants Coq (now Rocq), Isabelle, HOL Light and Lean have verified the four-colour theorem (Georges Gonthier, 2005), the Feit--Thompson odd order theorem (Gonthier and colleagues, 2012), the Kepler conjecture (Thomas Hales and colleagues, 2014), and the value of the fifth busy beaver (2024).

\begin{beauty}
Three things that look unrelated are one. A \emph{function} that takes a function and returns a function; a \emph{program} that runs, duplicates its arguments and throws some away; a \emph{proof} that assumes, discharges and concludes. Church's notation, invented to found logic on functions, failed as a logic and succeeded as a programming language, and then, with types, came back to logic as the language of proofs: modus ponens is function application, and simplifying a proof is running a program. The Church--Rosser theorem says that a computation's result does not depend on the order of its steps; Tait's theorem says that a typed computation always ends, which, read as logic, says that every proof can be made direct. The diagonal argument that gave Turing the halting problem gives here the $\mathsf Y$ combinator, recursion out of nothing.
\end{beauty}

\section{Exercises}

\begin{exercise}[$\star$]\label{ex:booleans}
Define $\mathsf{not}$, $\mathsf{and}$ and $\mathsf{or}$ on $\mathsf{true}=\lambda xy.x$ and $\mathsf{false}=\lambda xy.y$, and check them on all arguments.
\end{exercise}

\begin{exercise}[$\star$]\label{ex:exp}
Show that $\overline n\,\overline m\twoheadrightarrow\overline{m^n}$ for $n\ge1$ and all $m$.
\end{exercise}

\begin{exercise}[$\star$]\label{ex:diamond}
Show that one-step reduction does not have the diamond property: find $M\to N_1$ and $M\to N_2$ such that no $P$ satisfies $N_1\to P$ and $N_2\to P$.
\end{exercise}

\begin{exercise}[$\star$]\label{ex:turingfix}
Let $\Theta=AA$ with $A=\lambda xy.y(xxy)$. Show that $\Theta F\twoheadrightarrow F(\Theta F)$ for every $F$, and that $\mathsf YF\twoheadrightarrow F(\mathsf YF)$ fails for $F$ a variable.
\end{exercise}

\begin{exercise}[$\star\star$]\label{ex:minus}
Let $\mathsf{sub}=\lambda mn.n\,\mathsf{pred}\,m$. Show that $\mathsf{sub}\,\overline m\,\overline n\twoheadrightarrow\overline{m\dot-n}$, where $m\dot-n=\max(m-n,0)$, and $\lambda$-define the test $m\le n$ and equality of numbers.
\end{exercise}

\begin{exercise}[$\star\star$]\label{ex:strategies}
Show that $(\lambda x.y)\Omega$ has a normal form and an infinite reduction sequence, and that the leftmost strategy reaches the normal form of $\mathsf{K}\,\mathsf I\,\Omega$ and of $(\lambda x.y)\Omega$. Find a term with a normal form on which the strategy ``contract the rightmost innermost redex'' runs forever.
\end{exercise}

\begin{exercise}[$\star$]\label{ex:untypable}
Show that $\lambda x.xx$ is not typable, and deduce that $\Omega$ and $\mathsf Y$ are not typable.
\end{exercise}

\begin{exercise}[$\star\star$]\label{ex:types}
Show that $\vdash\mathsf S:(A\to B\to C)\to(A\to B)\to A\to C$ and that $\mathsf S\mathsf K\mathsf K\twoheadrightarrow\mathsf I$. Read the type of $\mathsf S$ as a formula: which tautology does $\mathsf S$ prove? Show that $\mathsf{plus}$ and $\mathsf{mult}$ have type $N_A\to N_A\to N_A$.
\end{exercise}

\begin{exercise}[$\star\star$]\label{ex:peirce}
Show that no closed term has type $((p\to q)\to p)\to p$, where $p\ne q$ are atomic, so that Peirce's law is not provable in minimal logic. Find a closed term of type $((p\to q)\to p)\to(p\to q)\to q$.
\end{exercise}

\begin{exercise}[$\star\star\star$]\label{ex:turingwn}
(Turing.) Prove directly that every typable term has a normal form. Define the \emph{degree} of a type by $\delta(p)=1$ and $\delta(A\to B)=1+\max(\delta(A),\delta(B))$, and the degree of a redex $(\lambda x.M)N$ in a typed term as the degree of the type of $\lambda x.M$. Show that contracting a redex of maximal degree $d$ that contains no other redex of degree $d$ in its argument $N$ creates only redexes of degree less than $d$, and conclude by induction on the pair (largest degree, number of redexes of that degree).
\end{exercise}

\section{Hints}

\begin{hint}{ex:booleans}
A boolean chooses: $\mathsf{and}\,b\,c$ is ``if $b$ then $c$ else false''.
\end{hint}
\begin{hint}{ex:exp}
Show by induction on $n$ that $\overline m^{\,n}(f)$, the $n$-fold application of $\overline m$ to $f$, reduces to $\lambda x.f^{m^n}(x)$.
\end{hint}
\begin{hint}{ex:diamond}
Look at Figure~\ref{fig:lambdagraph}.
\end{hint}
\begin{hint}{ex:turingfix}
Reduce the outer redex twice. For $\mathsf Y$, list all reducts of $\mathsf YF$.
\end{hint}
\begin{hint}{ex:minus}
$\overline n\,\mathsf{pred}\,\overline m$ applies $\mathsf{pred}$ $n$ times. Then $m\le n$ if and only if $m\dot-n=0$.
\end{hint}
\begin{hint}{ex:strategies}
For the last question, take a term which throws away an argument that has no normal form.
\end{hint}
\begin{hint}{ex:untypable}
Compare the sizes of $A$ and $A\to B$.
\end{hint}
\begin{hint}{ex:types}
Build the derivation from the bottom: $x:A\to B\to C$, $y:A\to B$, $z:A$.
\end{hint}
\begin{hint}{ex:peirce}
By Theorem~\ref{thm:sn} and Lemma~\ref{lem:normalshape}, look for a normal term; its shape is forced, and at some point a term of type $q$ is needed.
\end{hint}
\begin{hint}{ex:turingwn}
New redexes appear in two ways: when the argument $N$, an abstraction, is substituted for a variable in function position, and when the contractum, an abstraction, is applied to something.
\end{hint}

\section{Solutions}

\begin{solution}{ex:booleans}
\idea A boolean is a choice between two arguments, so each connective is a choice.
\writeup Let $\mathsf{not}=\lambda b.b\,\mathsf{false}\,\mathsf{true}$, $\mathsf{and}=\lambda bc.b\,c\,\mathsf{false}$ and $\mathsf{or}=\lambda bc.b\,\mathsf{true}\,c$. Then $\mathsf{not}\,\mathsf{true}\twoheadrightarrow\mathsf{true}\,\mathsf{false}\,\mathsf{true}\twoheadrightarrow\mathsf{false}$ and $\mathsf{not}\,\mathsf{false}\twoheadrightarrow\mathsf{true}$; $\mathsf{and}\,\mathsf{true}\,c\twoheadrightarrow c$ and $\mathsf{and}\,\mathsf{false}\,c\twoheadrightarrow\mathsf{false}$, which gives the four values; $\mathsf{or}\,\mathsf{true}\,c\twoheadrightarrow\mathsf{true}$ and $\mathsf{or}\,\mathsf{false}\,c\twoheadrightarrow c$.
\end{solution}

\begin{solution}{ex:exp}
\idea A numeral iterates; iterating ``iterate $m$ times'' $n$ times iterates $m^n$ times.
\writeup $\overline n\,\overline m\twoheadrightarrow\lambda x.\overline m^{\,n}(x)$, where $\overline m^{\,n}(x)$ is $\overline m(\overline m(\cdots(\overline mx)))$. We show by induction on $n\ge1$ that $\overline m^{\,n}(f)\twoheadrightarrow\lambda y.f^{m^n}(y)$ for a variable $f$. For $n=1$, $\overline mf\twoheadrightarrow\lambda y.f^m(y)$. If it holds for $n$, then $\overline m^{\,n+1}(f)=\overline m(\overline m^{\,n}(f))\twoheadrightarrow\overline m(\lambda y.f^{m^n}(y))\twoheadrightarrow\lambda y.(\lambda y.f^{m^n}(y))^m(y)\twoheadrightarrow\lambda y.f^{m^n\cdot m}(y)$. Hence $\overline n\,\overline m\twoheadrightarrow\lambda f.\overline m^{\,n}(f)\twoheadrightarrow\lambda fy.f^{m^n}(y)=\overline{m^n}$.
\end{solution}

\begin{solution}{ex:diamond}
\idea Duplicating a redex before contracting it costs two steps on the other side.
\writeup Take $M=(\lambda x.xx)(\mathsf I\mathsf I)$, $N_1=(\mathsf I\mathsf I)(\mathsf I\mathsf I)$ and $N_2=(\lambda x.xx)\mathsf I$. The only one-step reduct of $N_2$ is $\mathsf I\mathsf I$, and those of $N_1$ are $\mathsf I(\mathsf I\mathsf I)$ and $(\mathsf I\mathsf I)\mathsf I$, both different from $\mathsf I\mathsf I$. In two steps, both reach $\mathsf I\mathsf I$ (Figure~\ref{fig:lambdagraph}).
\end{solution}

\begin{solution}{ex:turingfix}
\idea $\Theta$ carries a copy of itself unevaluated, so that the fixed-point equation holds by reduction, not just by conversion.
\writeup $\Theta F=AAF\to(\lambda y.y(AAy))F\to F(AAF)=F(\Theta F)$. For $\mathsf Y$, let $W_f=(\lambda x.f(xx))(\lambda x.f(xx))$. The reducts of $\mathsf Y$ are the terms $Y_k=\lambda f.f^k(W_f)$, $k\ge0$ (with $Y_0=\mathsf Y$), and the reducts of $W_F$ are the terms $F^k(W_F)$. So the reducts of $\mathsf YF$ are the terms $Y_kF$ and $F^k(W_F)$: this set contains $\mathsf YF$ and is closed under one step, since the redexes of $Y_kF$ are the whole term, giving $F^k(W_F)$, and the one inside $W_f$. None of them is $F(\mathsf YF)$: the terms $Y_kF$ apply an abstraction, not $F$; and $F^k(W_F)=F(\mathsf YF)$ would need $k\ge1$ and $F^{k-1}(W_F)=\mathsf YF$, but $F^{k-1}(W_F)$ applies either $F$ or $\lambda x.F(xx)$, while $\mathsf YF$ applies $\mathsf Y$.
\end{solution}

\begin{solution}{ex:minus}
\idea Subtraction is iterated predecessor, and comparisons are tests on the difference.
\writeup $\mathsf{sub}\,\overline m\,\overline n\twoheadrightarrow\overline n\,\mathsf{pred}\,\overline m\twoheadrightarrow\mathsf{pred}^n(\overline m)$, and by Proposition~\ref{prop:pred} and induction on $n$, $\mathsf{pred}^n(\overline m)\twoheadrightarrow\overline{m\dot-n}$ (the predecessor of $\overline0$ is $\overline0$). Then $\mathsf{leq}=\lambda mn.\mathsf{iszero}(\mathsf{sub}\,m\,n)$ tests $m\le n$, and $\mathsf{eq}=\lambda mn.\mathsf{and}(\mathsf{leq}\,m\,n)(\mathsf{leq}\,n\,m)$ tests equality, with $\mathsf{and}$ from Exercise~\ref{ex:booleans}.
\end{solution}

\begin{solution}{ex:strategies}
\idea A redex whose argument is thrown away can be contracted first, and then the troublesome argument disappears.
\writeup $(\lambda x.y)\Omega\to y$ by the outer redex, which is normal, while contracting $\Omega$ again and again gives $(\lambda x.y)\Omega\to(\lambda x.y)\Omega\to\cdots$ forever. The leftmost redex of $(\lambda x.y)\Omega$ is the whole term, so the leftmost strategy reaches $y$ in one step; in $\mathsf K\mathsf I\Omega=(\mathsf K\mathsf I)\Omega$ the leftmost redex is $\mathsf K\mathsf I\to\lambda y.\mathsf I$, then $(\lambda y.\mathsf I)\Omega\to\mathsf I$. The rightmost innermost strategy contracts $\Omega$ first in $(\lambda x.y)\Omega$, since $\Omega$ is the only redex inside, and loops forever although the term has a normal form.
\end{solution}

\begin{solution}{ex:untypable}
\idea A type cannot be strictly contained in itself.
\writeup If $\Gamma\vdash\lambda x.xx:C$, then $C=A\to B$ and $\Gamma,x:A\vdash xx:B$, so $x$ has type $A$ and, being applied to $x$, also type $A\to B$, since the type of a variable is the one in the context. But $A\to B$ has more symbols than $A$, so $A\ne A\to B$. Every subterm of a typable term is typable (in an extended context), so $\Omega$ and $\mathsf Y$, which contain $\lambda x.xx$ or $\lambda x.f(xx)$ (same argument), are not typable.
\end{solution}

\begin{solution}{ex:types}
\idea Type checking is reading a proof backwards; the type of $\mathsf S$ is a tautology about implication.
\writeup In the context $x:A\to B\to C$, $y:A\to B$, $z:A$: $xz:B\to C$, $yz:B$, so $xz(yz):C$, and three abstractions give the type. Then $\mathsf S\mathsf K\mathsf K\twoheadrightarrow\lambda z.\mathsf Kz(\mathsf Kz)\twoheadrightarrow\lambda z.z=\mathsf I$. Read as a formula, $(A\to B\to C)\to(A\to B)\to A\to C$ is the distributivity of implication, the second axiom of Hilbert's calculus, and $\mathsf K$ proves the first, $A\to B\to A$. For $\mathsf{plus}=\lambda mnfx.mf(nfx)$, take $m,n:N_A$, $f:A\to A$, $x:A$: $nfx:A$, $mf:A\to A$, so $mf(nfx):A$ and $\mathsf{plus}:N_A\to N_A\to N_A$. For $\mathsf{mult}=\lambda mnf.m(nf)$: $nf:A\to A$, so $m(nf):A\to A$, and $\lambda f.m(nf):N_A$.
\end{solution}

\begin{solution}{ex:peirce}
\idea A normal proof of an implication starts by assuming its hypothesis; then the only way forward uses the hypothesis itself, and leads to a goal $q$ that nothing can produce.
\writeup Suppose $\vdash M:((p\to q)\to p)\to p$. By Theorem~\ref{thm:sn} and Proposition~\ref{prop:subjred} we may take $M$ normal; it is closed, so $M=\lambda f.P$ with $f:(p\to q)\to p\vdash P:p$, $P$ normal. As $p$ is atomic, $P$ is not an abstraction, so $P=yM_1\cdots M_k$ with $y$ a variable of the context, here $f$; its type forces $k=1$ and $P=fN$ with $f:(p\to q)\to p\vdash N:p\to q$. If $N=\lambda x.R$, then $R:q$ in the context $f:(p\to q)\to p$, $x:p$; if $N$ is not an abstraction, it is $fN'$ with type $p$, not $p\to q$. So we need a normal $R$ of atomic type $q$, which must be $yR_1\cdots R_j$ with $y\in\{f,x\}$; but the final result type of $f$ is $p$, and so is that of $x$, never $q$. Each case is impossible, so there is no such $M$. On the other hand $\lambda fg.g(f g)$ has type $((p\to q)\to p)\to(p\to q)\to q$: with $f:(p\to q)\to p$ and $g:p\to q$, $fg:p$ and $g(fg):q$.
\end{solution}

\begin{solution}{ex:turingwn}
\idea Contracting a redex can create new redexes and copy old ones, but the new ones have smaller degree, and if we choose well the copies do too; a lexicographic measure then decreases.
\writeup Types are preserved by reduction (Proposition~\ref{prop:subjred}), so degrees are well defined along the way. Let $d$ be the largest degree of a redex of the typed term $M$ and $n$ the number of redexes of degree $d$. Choose a redex $R=(\lambda x.P)N$ of degree $d$ such that no redex of degree $d$ occurs inside $N$; for instance the one whose $\lambda$ is rightmost among the redexes of degree $d$, since the redexes inside $N$ begin to the right of it. Let $\lambda x.P:A\to B$, so $d=\delta(A\to B)$. Contracting $R$ into $P[x:=N]$: the redexes of $M$ outside $R$ remain, with their degrees (one of them may have $R$ as a subterm, which changes but keeps its type); the redexes inside $P$ remain, once each; the redexes inside $N$ are copied as many times as $x$ occurs, but they have degree less than $d$. New redexes can appear in two ways: if $N$ is an abstraction and $P$ contains $xQ$, the new redex $NQ'$ has degree $\delta(A)<d$; if the contractum $P[x:=N]$ is an abstraction and $R$ was applied, $RQ$, the new redex $P[x:=N]\,Q$ has degree $\delta(B)<d$. So the number of redexes of degree $d$ decreases by one, and no redex of degree $\ge d$ is created. After $n$ such steps the largest degree is less than $d$; by induction on $(d,n)$ in lexicographic order, the process ends at a term without redexes, a normal form.
\end{solution}

\section*{Sources and further reading}

\begin{itemize}
\item \emph{The stories.} A.~Church, \emph{A set of postulates for the foundation of logic}, Annals of Mathematics 33 (1932), 346--366, and second paper, ibid.\ 34 (1933), 839--864; S.~C.~Kleene and J.~B.~Rosser, \emph{The inconsistency of certain formal logics}, Annals of Mathematics 36 (1935), 630--636; S.~C.~Kleene, \emph{Origins of recursive function theory}, Annals of the History of Computing 3 (1981), 52--67, for the predecessor and Church's thesis; M.~Davis, \emph{Why Gödel didn't have Church's thesis}, Information and Control 54 (1982), 3--24, which quotes Church's letter to Kleene of 29 November 1935 on Gödel's reaction; J.~McCarthy, \emph{Recursive functions of symbolic expressions and their computation by machine, Part I}, Communications of the ACM 3 (1960), 184--195; H.~P.~Barendregt, \emph{The impact of the lambda calculus in logic and computer science}, Bulletin of Symbolic Logic 3 (1997), 181--215.
\item A.~Church, \emph{A formulation of the simple theory of types}, Journal of Symbolic Logic 5 (1940), 56--68; A.~Church, \emph{An unsolvable problem of elementary number theory}, American Journal of Mathematics 58 (1936), 345--363; A.~Church and J.~B.~Rosser, \emph{Some properties of conversion}, Transactions of the American Mathematical Society 39 (1936), 472--482; S.~C.~Kleene, \emph{$\lambda$-definability and recursiveness}, Duke Mathematical Journal 2 (1936), 340--353; A.~M.~Turing, \emph{Computability and $\lambda$-definability}, Journal of Symbolic Logic 2 (1937), 153--163, and \emph{The $\mathfrak p$-function in $\lambda$-$K$-conversion}, ibid., 164.
\item M.~Takahashi, \emph{Parallel reductions in $\lambda$-calculus}, Information and Computation 118 (1995), 120--127, the proof of Theorem~\ref{thm:churchrosser}; H.~B.~Curry and R.~Feys, \emph{Combinatory Logic}, vol.~I, North-Holland, 1958, for the normalisation theorem; H.~P.~Barendregt, \emph{The Lambda Calculus: Its Syntax and Semantics}, North-Holland, revised edition 1984, the reference work.
\item W.~W.~Tait, \emph{Intensional interpretations of functionals of finite type I}, Journal of Symbolic Logic 32 (1967), 198--212; J.-Y.~Girard, Y.~Lafont and P.~Taylor, \emph{Proofs and Types}, Cambridge University Press, 1989, whose Chapter~6 we follow for Theorem~\ref{thm:sn} and whose Chapter~4 gives Turing's proof of Exercise~\ref{ex:turingwn}; R.~O.~Gandy, \emph{An early proof of normalization by A.~M.~Turing}, in \emph{To H.~B.~Curry: Essays on Combinatory Logic, Lambda Calculus and Formalism}, Academic Press, 1980, 453--455; H.~Schwichtenberg, \emph{Definierbare Funktionen im $\lambda$-Kalkül mit Typen}, Archiv für mathematische Logik 17 (1976), 113--114.
\item H.~B.~Curry, \emph{Functionality in combinatory logic}, Proceedings of the National Academy of Sciences 20 (1934), 584--590; W.~A.~Howard, \emph{The formulae-as-types notion of construction}, manuscript 1969, in \emph{To H.~B.~Curry} (above), 479--490; G.~Gentzen, \emph{Untersuchungen über das logische Schließen}, Mathematische Zeitschrift 39 (1935), 176--210 and 405--431; T.~G.~Griffin, \emph{A formulae-as-types notion of control}, Proceedings of the 17th ACM Symposium on Principles of Programming Languages (1990), 47--58; M.~H.~S{\o}rensen and P.~Urzyczyn, \emph{Lectures on the Curry--Howard Isomorphism}, Elsevier, 2006.
\item J.~R.~Hindley, \emph{The principal type-scheme of an object in combinatory logic}, Transactions of the American Mathematical Society 146 (1969), 29--60; R.~Milner, \emph{A theory of type polymorphism in programming}, Journal of Computer and System Sciences 17 (1978), 348--375.
\item N.~G.~de Bruijn, \emph{A survey of the project Automath}, in \emph{To H.~B.~Curry} (above), 579--606; L.~S.~van Benthem Jutting, \emph{Checking Landau's ``Grundlagen'' in the Automath system}, PhD thesis, Eindhoven, 1977; G.~Gonthier, \emph{Formal proof---the four-color theorem}, Notices of the American Mathematical Society 55 (2008), 1382--1393; G.~Gonthier et al., \emph{A machine-checked proof of the odd order theorem}, Interactive Theorem Proving, Lecture Notes in Computer Science 7998 (2013), 163--179; T.~Hales et al., \emph{A formal proof of the Kepler conjecture}, Forum of Mathematics, Pi 5 (2017), e2.
\end{itemize}
